\section{From non-density to maximal atypical finiteness}\label{sec:finiteness}

The sequence theorem and its reduction in Section~\ref{sec:sequences}
prove Theorem~\ref{thm:nondensity} for every abelian variety over $\Qbar$.
We now explain why this universal statement gives the finiteness
asserted in Theorem~\ref{thm:main}. We use the reduction of Barroero
and Dill \cite[Theorem~1.9]{BarroeroDill2021}; their passage from
optimal points to positive-dimensional optimal subvarieties uses
Habegger and Pila's \cite[Theorem~9.8(i)]{HabeggerPila2016}.

For an irreducible subvariety $Z$ of an abelian variety, its
\emph{special defect} is
\[
 \delta(Z)=\dim\langle Z\rangle_{\mathrm{sp}}-\dim Z.
\]
An irreducible $Z\subseteq X$ is \emph{optimal in $X$} if every
irreducible $U$ with $Z\subsetneq U\subseteq X$ satisfies
$\delta(U)>\delta(Z)$. This uses special closure, in contrast to the
arbitrary-coset closure used for geodesic optimality in
Section~\ref{sec:tangent}.

\begin{proposition}\label{prop:optimal-finite}
Every irreducible subvariety of every abelian variety over $\Qbar$
contains only finitely many optimal subvarieties.
\end{proposition}

\begin{proof}
Fix an abelian variety $A$, and nonnegative integers $m,e$.
Theorem~1.9 of \citet{BarroeroDill2021} has the following implication.
Suppose that, for every abelian subvariety $B\subseteq A$ and every
irreducible $Z\subseteq B$ of dimension at most $m$ which is not
contained in a proper special subvariety of $B$, the set
\[
 Z\cap B^{[\max\{\dim Z+1,\,\dim B-e\}]}
\]
is not dense in $Z$. Then each irreducible subvariety of $A$ of
dimension at most $m$ has only finitely many optimal subvarieties of
defect at most $e$.

The displayed intersection is a subset of
$Z\cap B^{[\dim Z+1]}$, so Theorem~\ref{thm:nondensity} verifies the
hypothesis for every $B$ and $Z$ in this statement. All these varieties
are defined over number fields after finite extension. The same
non-density theorem is available for quotients and isogenous varieties
as well, so no ambient restriction is introduced by the reduction.
Take $m=\dim X$ and $e=\dim A$. Every special defect in $A$ lies between
$0$ and $\dim A$, which proves the proposition.
\end{proof}

\begin{lemma}\label{lem:maximal-optimal}
Let $X$ be an irreducible subvariety of an abelian variety and
$S=\langle X\rangle_{\mathrm{sp}}$. Every maximal atypical subvariety
of $X$ in $S$ is optimal in $X$.
\end{lemma}

\begin{proof}
Let $Y$ be maximal atypical, witnessed by a special $T\subseteq S$.
Since $\langle Y\rangle_{\mathrm{sp}}\subseteq T$, inequality
\eqref{eq:atypical} gives
\[
 \delta(Y)\le\dim T-\dim Y
 <\dim S-\dim X=\delta(X).
\]
Suppose $Y\subsetneq U\subseteq X$ and
$\delta(U)\le\delta(Y)$. Put $T'=\langle U\rangle_{\mathrm{sp}}$;
then $T'\subseteq S$ by \eqref{eq:closure-monotonicity}.
Choose an irreducible component $C$ of $X\cap T'$ containing $U$.
We have
\[
 \dim C\ge\dim U
 =\dim T'-\delta(U)
 >\dim T'-\delta(X)
 =\dim X+\dim T'-\dim S.
\]
Thus $C$ is atypical and strictly contains $Y$, a contradiction.
\end{proof}

\begin{proof}[Proof of Theorem~\ref{thm:main}]
Write $S=B+\tau$ with $\tau$ torsion. Translation by $-\tau$ identifies
$S$ with the abelian variety $B$, preserves dimensions and special
subvarieties, and sends $X$ to a subvariety with special closure $B$.
The torsion point, $B$ and this translated subvariety are all defined
over a number field after a finite extension. We may therefore work
inside $B$.

Proposition~\ref{prop:optimal-finite} and
Lemma~\ref{lem:maximal-optimal} show that the set of maximal atypical
subvarieties is finite. Each is proper: if $X$ were a component
witnessing its own atypicality, the witnessing special subvariety
would contain $X$ and hence equal its special closure, making the
strict dimension inequality impossible. Starting from any atypical
subvariety, a sequence of strict atypical enlargements increases
dimension and must terminate. Thus every atypical subvariety lies in a
maximal one. Their finite union is a proper closed subset of the
irreducible variety $X$.

If $\dim X=0$, a special subvariety containing its only point must
contain its special closure, so there are no atypical subvarieties.
This also covers the zero-dimensional ambient case. The assertions
and the possibility of an empty exceptional list follow in every case.
\end{proof}
